\subsection{Rúbricas como señales de recompensa}
Una línea de investigación reciente busca extender el aprendizaje por refuerzo a tareas abiertas mediante recompensas basadas en rúbricas, que descomponen la calidad de una respuesta en criterios explícitos. Dentro de esta línea, Huang y sus colaboradores presentaron Rubicon, un enfoque que utiliza rúbricas elaboradas por personas, generadas por modelos de lenguaje o construidas mediante una combinación de ambos para orientar el entrenamiento en tareas de escritura, humanidades y seguimiento de instrucciones. Su propuesta combina criterios generales con otros específicos de cada tarea o consulta, permitiendo evaluar propiedades que no pueden reducirse a la comprobación de una única respuesta correcta. Para ello, construyeron un banco de más de 10.000 rúbricas y entrenaron Qwen3-30B-A3B con algo más de 5.000 ejemplos, obteniendo un aumento del puntaje promedio de 65,29 a 70,50 en las siete evaluaciones abiertas reportadas, aunque el puntaje en IFEval descendió de 83,55 a 81,70 \cite{huang2025rubicon}. Estos resultados muestran mejoras en el promedio junto con retrocesos en una evaluación particular, y los autores señalan que incorporar más rúbricas no garantiza por sí solo mejores resultados, ya que también importan su diversidad, su nivel de detalle y la estrategia utilizada durante el entrenamiento. Este antecedente resulta relevante para el presente trabajo porque muestra cómo las rúbricas pueden emplearse como señales de optimización, aunque su aplicación a consultas clínicas requiere considerar criterios y riesgos propios del dominio.

